\documentclass{article}

% A synthetic seed paper for Lattice's research environment. Its numbers are
% illustrative, not measured results.
\usepackage[preprint]{neurips_2026}
\usepackage[utf8]{inputenc}
\usepackage[T1]{fontenc}
\usepackage{amsmath,amssymb}
\usepackage{amsthm}
\usepackage{booktabs}
\usepackage{multirow}
\usepackage{graphicx}
\usepackage{algorithm}
\usepackage{algpseudocode}
\usepackage{xcolor}
\usepackage{hyperref}
\usepackage{cleveref}

\newtheorem{proposition}{Proposition}
\newtheorem{lemma}{Lemma}
\DeclareMathOperator*{\argmax}{arg\,max}
\DeclareMathOperator{\topk}{top\text{-}k}
\newcommand{\R}{\mathbb{R}}
\newcommand{\expert}[1]{E_{#1}}

\title{Sparse Routing for Efficient Sequence Models}
\author{Seed Author \\
Lattice Research Environment \\
\texttt{seed@lattice.example}}

\begin{document}
\maketitle

\begin{abstract}
Mixture-of-experts layers route each token to a small subset of expert networks, so model capacity can grow without a matching growth in computation per token.
This note studies a simple balanced router that combines top-$k$ selection with a capacity limit per expert.
We describe the router, prove that its load stays within a constant factor of the uniform load, and compare it with dense layers of equal cost on two synthetic sequence tasks.
\end{abstract}

\section{Introduction}
\label{sec:intro}

Conditional computation activates only part of a network for each input.
Sparsely gated mixture-of-experts layers made this practical at scale by learning a gate that picks a few experts per token \citep{shazeer2017outrageously}.
Later systems simplified routing to a single expert per token \citep{fedus2022switch} or sharded experts across many accelerators \citep{lepikhin2021gshard}.
In each case the router must avoid sending most tokens to a few popular experts, a failure that wastes capacity and stalls training.

Our contributions are:
\begin{itemize}
  \item a router that pairs top-$k$ selection with a hard capacity per expert (\cref{sec:method});
  \item a bound on the maximum expert load under that router (\cref{prop:load});
  \item a comparison with dense layers of equal floating-point cost (\cref{sec:experiments}).
\end{itemize}

\section{Balanced top-$k$ routing}
\label{sec:method}

Let $x_t \in \R^d$ be the representation of token $t$ and let $W_r \in \R^{d \times n}$ be the router's weights for $n$ experts.
The router scores the experts with
\begin{equation}
  s_t = \operatorname{softmax}(W_r^\top x_t),
  \label{eq:scores}
\end{equation}
and the layer's output is the weighted sum of the selected experts,
\begin{align}
  y_t &= \sum_{i \in S_t} s_{t,i}\, \expert{i}(x_t), \label{eq:output}\\
  S_t &= \topk_{i \le n} \, s_{t,i}. \nonumber
\end{align}
Each expert $\expert{i}$ is a two-layer feed-forward network, as in the Transformer block.

\paragraph{Capacity.}
A batch of $T$ tokens gives every expert a capacity of $C = \lceil c\,kT/n \rceil$ slots for a capacity factor $c \ge 1$.
Tokens are assigned in order of their score; once an expert is full, a token falls back to its next-best expert, and a token with no free expert in its top $k$ is passed through the residual connection unchanged.
\Cref{alg:router} lists the procedure.

\begin{algorithm}[t]
  \caption{Balanced top-$k$ routing for one batch.}
  \label{alg:router}
  \begin{algorithmic}[1]
    \Require scores $s_t$ from \cref{eq:scores}, capacity $C$
    \State $\mathrm{load}_i \gets 0$ for every expert $i$
    \For{each token $t$ in decreasing order of $\max_i s_{t,i}$}
      \For{each expert $i$ in $\topk_i s_{t,i}$}
        \If{$\mathrm{load}_i < C$}
          \State assign $t$ to $i$ and set $\mathrm{load}_i \gets \mathrm{load}_i + 1$
        \EndIf
      \EndFor
    \EndFor
  \end{algorithmic}
\end{algorithm}

\begin{proposition}
\label{prop:load}
Under balanced top-$k$ routing with capacity factor $c$, no expert receives more than $c\,kT/n + 1$ tokens in a batch of $T$ tokens.
\end{proposition}

\begin{proof}
An expert accepts a token only while its load is below $C$, so its final load is at most $C \le c\,kT/n + 1$.
\end{proof}

The bound holds for any input, which is what lets the router drop the auxiliary loss used by \citet{fedus2022switch}.
\Cref{lem:dropped} bounds the share of tokens that pass through unrouted.

\begin{lemma}
\label{lem:dropped}
If every token's top $k$ experts are drawn uniformly at random, the expected share of unrouted tokens decreases monotonically in $c$.
\end{lemma}

\section{Experiments}
\label{sec:experiments}

We compare the balanced router with a dense feed-forward layer of equal cost and with unconstrained top-$k$ routing.\footnote{All numbers in this seed paper are illustrative and were not measured.}
Both tasks are synthetic: copying a sequence with long gaps, and predicting the parity of a sliding window.
Every model has four layers and model width $d = 256$; the mixture models use $n = 16$ experts with $k = 2$.

\begin{figure}[t]
  \centering
  \includegraphics[width=0.8\linewidth]{figures/routing.png}
  \caption{Expert load over training for unconstrained and balanced routing. The balanced router stays near the uniform load.}
  \label{fig:load}
\end{figure}

\Cref{fig:load} shows the load of the busiest expert during training.
Without a capacity limit, two experts receive most tokens within the first thousand steps; with $c = 1.25$ the busiest expert stays within the bound of \cref{prop:load}.
\Cref{tab:results} reports accuracy at the end of training.

\begin{table}[t]
  \centering
  \caption{Accuracy (\%) on the two synthetic tasks, mean of three runs (illustrative).}
  \label{tab:results}
  \begin{tabular}{@{}llcc@{}}
    \toprule
    \multicolumn{2}{@{}l}{Model} & Copy & Parity \\
    \midrule
    Dense & $d_\mathrm{ff} = 1024$ & 91.2 & 74.5 \\
    \multirow{2}{*}{Top-$k$} & unconstrained & 93.8 & 77.1 \\
     & balanced, $c = 1.25$ & 95.4 & 79.6 \\
    \bottomrule
  \end{tabular}
\end{table}

The balanced router improves on both baselines, and the gain is largest on the parity task, where unconstrained routing leaves most experts idle.


\section{Related work}
\label{sec:related}

The idea of combining specialised experts with a learned gate predates deep learning \citep{jacobs1991adaptive}.
Transformer-based mixture models inherit the attention layer of \citet{vaswani2017attention} and replace its feed-forward block with experts.
Load balancing has been handled with auxiliary losses \citep{shazeer2017outrageously,fedus2022switch}, with capacity factors \citep{lepikhin2021gshard}, and with routing formulated as an assignment problem \citep{lewis2021base}.
Our router follows the capacity approach but analyses its worst-case load directly.

\section{Conclusion}
\label{sec:conclusion}

A capacity limit turns top-$k$ routing into a procedure whose load is bounded by construction, which removes the need to tune an auxiliary balancing loss.
The bound in \cref{prop:load} is loose for skewed inputs, and tightening it is left to future work.

\bibliographystyle{plainnat}
\bibliography{references}

\appendix
\section{Proof details}
\label{app:proof}

The capacity rule of \cref{sec:method} assigns a token to an expert only when that expert has a free slot.
Because $C = \lceil c\,kT/n \rceil$, we have $C \le c\,kT/n + 1$, which gives \cref{prop:load}.
For \cref{lem:dropped}, a token is unrouted only if all of its $k$ choices are full, and the probability of that event falls as each capacity $C$ grows with $c$.


\end{document}
